\begin{abstract}
Training-free open-vocabulary semantic segmentation must assign arbitrary textual categories to dense remote-sensing imagery without updating the segmentation model. Concept-grounded segmenters such as SAM3 and SegEarth-OV3 provide a shared grounding state with distinct Presence, Semantic, and Instance readouts, but current methods typically condition all readouts with one category-level description and one visual observation. We study the resulting mismatch between shared concept grounding and heterogeneous evidence requirements. We introduce Joint Role--View Specialization (JRV), which preserves the official class anchor, grounds role-functional candidates with the complete frozen model, and applies bounded role-specific corrections while retaining the native Presence--Instance coupling. JRV then grounds spatially aligned Local and Context observations, reuses the same role profile across both views, and composes their complete class-score fields before the original class decision. On eight remote-sensing benchmarks, JRV raises the matched average mIoU from 53.56 to 57.49 and improves every dataset over the baseline. A matched Role--View factorial analysis shows positive interaction on six of eight datasets, while the natural-image evaluations show transfer of role supplementation and the dataset dependence of view composition. These results support organizing supplementary evidence by readout responsibility and observation condition rather than treating additional text or larger visual fields as uniformly beneficial. Our code is available at \url{https://github.com/pengjunhao621-arch/SegEarth3}.
\end{abstract}
